% ##############################################################################
\section{Related work}
\label{sec:related_work}

% ==============================================================================
\subsection{LLM gateways and routing}
\label{sec:related_gateways}

Multi-provider gateways such as OpenRouter unify access to many LLM
providers behind a single API, which is the starting point for \NoesisServer{}
(Section~\ref{sec:gateway_architecture}); unlike \Noesis, they typically set
prices unilaterally rather than clearing an auction, and do not meter delivered
quality against a contract. RouteLLM \cite{routellm2024} and FrugalGPT
\cite{frugalgpt2023} address the difficulty-aware routing problem of
Section~\ref{sec:routing} directly, learning to route or cascade requests
across models of different cost to reduce spend while controlling quality loss;
\Noesis{} adopts a similar routing formalization but ties the resulting cost
savings to contract fulfillment and to a market that prices the tiers being
routed between.

% ==============================================================================
\subsection{Model fusion and distillation}
\label{sec:related_fusion}

Self-consistency \cite{wang2022selfconsistency} and mixture-of-agents
\cite{wang2024moa} are examples of the answer-fusion strategy of
Section~\ref{sec:fusion}: combining multiple model outputs, by voting or by
a further aggregation model, to trade additional inference cost for
higher answer quality. Knowledge distillation \cite{hinton2015} underlies
the training-signal reuse of Section~\ref{sec:distillation}, in which
\NoesisServer{}'s logged traffic plays the role of the training corpus.

% ==============================================================================
\subsection{Auction and market microstructure}
\label{sec:related_auctions}

The uniform-price call auction of Problem~\ref{prob:clearing} is a
multi-attribute generalization of the classical double auction
\cite{mcafee1992}. The choice of a periodic batch auction, rather than a
continuous limit order book, follows the argument that frequent batch
auctions remove the latency-arbitrage incentives that plague continuous
markets \cite{budish2015}, an argument we expect to carry over to a
market where sellers could otherwise race to snipe favorably priced
demand. The strategic risks flagged in
Section~\ref{sec:mechanism_design_risks}, bid shading and collusion, draw
on the classical auction-theory literature on optimal and
incentive-compatible auction design \cite{myerson1981} and on collusive
bidding \cite{mcafeemcmillan1992}.

% ==============================================================================
\subsection{Commodity and capacity market design}
\label{sec:related_commodity}

Electricity market design \cite{stoft2002, cramton2017} is the closest
analogy to \Noesis, as discussed in
Section~\ref{sec:related_commodity_markets}: the day-ahead/real-time
separation motivates the tiered, periodic clearing of
Section~\ref{sec:rounds}, and non-performance penalties in capacity
markets \cite{cramtonstoft2005} motivate the reputation and pricing
feedback loop of Section~\ref{sec:reputation}.

% ==============================================================================
\subsection{Decentralized exchange protocols}
\label{sec:related_dex}

\Noesis{}'s contract, bid, and ask notation (Section~\ref{sec:contracts_and_notation})
and its treatment of the clearing problem as a linear program
(Problem~\ref{prob:clearing}) follow the approach taken by Tulip, a
protocol for decentralized token exchange built around limit orders and a
Walrasian-style auction \cite{tulip2023}. The two protocols differ in an
important respect: Tulip's traded good, a token, is settled entirely
on-chain and requires no external verification, whereas \Noesis's traded
good, a unit of intelligence, is produced by off-chain computation whose
quality can only be estimated, not settled atomically, which is why
\NoesisServer's fulfillment monitoring (Section~\ref{sec:fulfillment_monitoring})
has no counterpart in Tulip. Whether \NoesisMarket's contract matching and
settlement, as opposed to the inference computation itself, should move
on-chain, in the manner of Tulip or of automated market makers such as
Uniswap \cite{uniswap2020, angerischitra2020}, is sketched as a concrete,
blind-bid, stake-backed design in Section~\ref{sec:decentralized_extension},
drawing on the companion ZK-Tulip design notes \cite{zktulip2023}; the
resulting open questions are consolidated in
Section~\ref{sec:open_questions}.

Figure~\ref{fig:related_landscape} places the work surveyed above on
these two axes: routing/fusion work and auction theory each occupy one
axis; electricity and capacity markets are the closest existing analogy,
occupying both, but for a physically metered commodity rather than
machine intelligence (Section~\ref{sec:related_commodity}); \Noesis{} is
the only point combining both axes for machine intelligence.

% rendered_images:begin
% ```tikz
% \definecolor{axiscolor}{RGB}{133,146,163}
% \definecolor{dividercolor}{RGB}{199,208,218}
% \definecolor{ptcolor}{RGB}{31,78,121}
% \definecolor{noesiscolor}{RGB}{192,69,91}
% \fontfamily{phv}\selectfont
% \draw[->, >=stealth, thick, axiscolor] (0,0) -- (7.2,0) node[right, font=\small, align=left] {Auction-based\\price discovery};
% \draw[->, >=stealth, thick, axiscolor] (0,0) -- (0,5.55) node[above, font=\small, align=center] {Metered, verifiable\\fulfillment};
% \draw[dashed, dividercolor] (3.45,-0.225) -- (3.45,5.55);
% \draw[dashed, dividercolor] (-0.225,2.7) -- (7.2,2.7);
% \node[font=\small] at (1.725,-0.375) {No};
% \node[font=\small] at (5.25,-0.375) {Yes};
% \node[font=\small] at (-0.45,1.35) {No};
% \node[font=\small] at (-0.45,4.2) {Yes};
% \filldraw[ptcolor] (1.275,0.75) circle (2.4pt);
% \node[anchor=west, align=left, font=\small] at (0.225,1.425) {LLM routing \& fusion\\(RouteLLM, FrugalGPT,\\self-consistency, MoA)};
% \filldraw[ptcolor] (5.475,0.75) circle (2.4pt);
% \node[anchor=west, align=left, font=\small] at (4.725,1.425) {Auction theory\\(McAfee, Budish,\\Myerson)};
% \filldraw[ptcolor] (2.925,0.75) circle (2.4pt);
% \node[anchor=west, align=left, font=\small] at (2.175,0.225) {Decentralized exchange\\(Tulip, Uniswap)};
% \filldraw[ptcolor] (5.4,4.05) circle (2.4pt);
% \node[anchor=east, align=right, font=\small] at (4.95,4.5) {Commodity \& capacity\\markets (electricity)};
% \filldraw[noesiscolor] (6.6,4.95) circle (3pt);
% \node[anchor=west, align=left, font=\small\bfseries, text=noesiscolor] at (6.225,5.29) {Noesis};
% ```
% label=fig:related_landscape
% caption=Related work positioned by auction-based price discovery and metered fulfillment.
% rendered_images:end
% render_images:begin
\begin{figure}[H]
  \includegraphics[width=\linewidth]{06_related_work.tex.figs/06_related_work.1.png}
  \caption{Related work positioned by auction-based price discovery and metered fulfillment.}
  \label{fig:related_landscape}
\end{figure}
% render_images:end
